\documentclass[12pt,a4paper]{article}

\usepackage[margin=1in]{geometry}
\usepackage[T1]{fontenc}
\usepackage[utf8]{inputenc}
\usepackage{lmodern}
\usepackage{setspace}
\usepackage{graphicx}
\usepackage{enumitem}
\usepackage{array}
\usepackage{longtable}
\usepackage{booktabs}
\usepackage{titlesec}
\usepackage[hidelinks]{hyperref}

\setstretch{1.15}
\setlength{\parindent}{0pt}
\setlength{\parskip}{6pt}

\titleformat{\section}{\normalfont\large\bfseries}{\thesection}{1em}{}
\titleformat{\subsection}{\normalfont\normalsize\bfseries}{\thesubsection}{1em}{}

% ---------- diagram helper ----------
% Usage: \insertdiagram{filename.png}{Caption text}
\newcommand{\insertdiagram}[2]{%
    \begin{center}
    \IfFileExists{#1}{
        \includegraphics[width=0.9\textwidth]{#1}
    }{
        \fbox{%
            \parbox[c][6cm][c]{0.88\textwidth}{\centering
            Diagram placeholder\\[6pt]
            Insert PNG file here: \texttt{#1}}
        }
    }
    \end{center}
    \begin{center}
    \textit{#2}
    \end{center}
}

\begin{document}

% ================= TITLE PAGE =================
\begin{center}
\textbf{UNIVERSITY OF CHITTAGONG}\\
Department of Computer Science and Engineering
\end{center}

\vspace{1cm}

\begin{center}
\textbf{\Large Assignment}
\end{center}

\vspace{0.8cm}

\textbf{Course Title :} Database Management System\\
\textbf{Course Code No. :} CSE 414

\vspace{1cm}

\textbf{Submitted to:}\\
Dr. Rudra Pratap Deb Nath\\
Professor\\
Department of Computer Science and Engineering\\
University of Chittagong

\vspace{1cm}

\textbf{Submitted by:}\\
Name: Sheikh Tawsif Bin Aftab\\
ID: 24701026\\
Session: 2023--2024 B.Sc. Engineering\\
(4th Semester)\\
Department of Computer Science and Engineering\\
University of Chittagong

\vspace{1cm}

\textbf{Date of Submission:}\\
28 February 2026

\newpage

% ================= MAIN BODY =================
\begin{center}
\textbf{\Large Chapter 6}
\end{center}

\textbf{Question 1}\\
Analyze, solve, and present the solution of Problem 6.1.

\textbf{Answer:}

For the car insurance company, the main entity sets are \textbf{Customer}, \textbf{Car}, \textbf{Accident}, \textbf{Policy}, and \textbf{Payment}.

\insertdiagram{q6_1.png}{E-R diagram for Problem 6.1}

\textbf{Entities:}
\begin{itemize}[leftmargin=1.5em]
    \item Customer (\underline{customer\_id}, name, address, phone)
    \item Car (\underline{car\_id}, license\_no, make, model, year)
    \item Accident (\underline{accident\_id}, accident\_date, location, description)
    \item Policy (\underline{policy\_no}, issue\_date, type, status)
    \item Payment (\underline{payment\_id}, period\_from, period\_to, due\_date, received\_date, amount)
\end{itemize}

\textbf{Relationships:}
\begin{itemize}[leftmargin=1.5em]
    \item Owns(Customer, Car): one customer can own many cars
    \item Has\_Accident(Car, Accident): one car can have many accidents
    \item Covers(Policy, Car): many-to-many
    \item Has\_Payment(Policy, Payment): one policy has many payments
\end{itemize}

\newpage
\textbf{Question 2}\\
Analyze, solve, and present the solution of Problem 6.2.

\textbf{Answer:}

\textbf{(a) Using a ternary relationship with Exam as an entity}

\insertdiagram{q6_2a.png}{E-R diagram for Problem 6.2(a)}

The required entity sets are:
\begin{itemize}[leftmargin=1.5em]
    \item Student (\underline{student\_id}, name)
    \item Course (\underline{course\_id}, title)
    \item Section (\underline{section\_id}, sec\_id, semester, year)
    \item Exam (\underline{exam\_id}, exam\_type, exam\_date, max\_marks)
\end{itemize}

A ternary relationship \textbf{Marks(Student, Section, Exam)} is used with attribute \textbf{score}.  
This is appropriate because a student's mark depends on the student, the section, and the exam together.

\textbf{(b) Using only a binary relationship between Student and Section}

\insertdiagram{q6_2b.png}{E-R diagram for Problem 6.2(b)}

A suitable design is to introduce:
\begin{itemize}[leftmargin=1.5em]
    \item Enrollment (\underline{enrollment\_id}, student\_id, section\_id)
    \item Exam (\underline{exam\_id}, section\_id, exam\_type, exam\_date, max\_marks)
    \item Exam\_Result (enrollment\_id, exam\_id, score)
\end{itemize}

This design keeps one binary student-section relationship while allowing multiple exam results.

\textbf{Question 3}\\
Analyze, solve, and present the solution of Problem 6.3.

\textbf{Answer:}

\insertdiagram{q6_3.png}{E-R diagram for Problem 6.3}

For the sports team scoring system, the main entity sets are:
\begin{itemize}[leftmargin=1.5em]
    \item Team (\underline{team\_id}, name, city)
    \item Player (\underline{player\_id}, name, position, jersey\_no)
    \item Match (\underline{match\_id}, match\_date, venue, opponent, team\_score, opponent\_score)
\end{itemize}

A relationship \textbf{Participation(Player, Match)} stores match-specific statistics such as:
\begin{itemize}[leftmargin=1.5em]
    \item minutes\_played
    \item started\_flag
    \item goals
    \item assists
    \item penalty\_goals
    \item own\_goals
\end{itemize}

Derived values such as total goals, total assists, and total wins should be computed from the stored data.

\textbf{Question 4}\\
Analyze, solve, and present the solution of Problem 6.5.

\textbf{Answer:}

\textbf{(a) Disconnected graph:}  
A disconnected E-R graph means that two or more parts of the design are not linked by any relationship. This may indicate either independent subdomains or a missing relationship.

\textbf{(b) Cycle in a graph:}  
A cycle means that by traversing relationships, it is possible to return to the same entity set. This is not necessarily an error; it often reflects valid semantic connections.

\textbf{Question 5}\\
Analyze, solve, and present the solution of Problem 6.7.

\textbf{Answer:}

A weak entity should not always be converted into a strong entity because doing so introduces redundancy and hides existence dependence.

\textbf{Example:}
\begin{itemize}[leftmargin=1.5em]
    \item Employee (\underline{emp\_id}, name)
    \item Dependent is a weak entity identified through Employee
\end{itemize}

If Dependent is converted into a strong entity as  
Dependent(emp\_id, dependent\_name, age),  
then emp\_id must be repeated for every dependent.

Therefore, weak entities are useful because they clearly represent:
\begin{itemize}[leftmargin=1.5em]
    \item existence dependence
    \item identification through the owner
    \item cleaner conceptual modeling
\end{itemize}

\textbf{Question 6}\\
Analyze, solve, and present the solution of Problem 6.9.

\textbf{Answer:}

Suppose the relation is:
advisor(s\_ID, i\_ID)

To make the advisor relationship one-to-one:
\begin{itemize}[leftmargin=1.5em]
    \item each student must have at most one instructor
    \item each instructor must advise at most one student
\end{itemize}

So:
\begin{itemize}[leftmargin=1.5em]
    \item s\_ID must be unique
    \item i\_ID must also be unique
\end{itemize}

A suitable relational form is:
advisor(s\_ID PRIMARY KEY, i\_ID UNIQUE)

\newpage
\textbf{Question 7}\\
Analyze, solve, and present the solution of Problem 6.13*.

\textbf{Answer:}

\insertdiagram{q6_13.png}{Temporal E-R design for Problem 6.13}

To represent temporal information, valid time should be associated with entities, attributes, and relationship sets.

\textbf{Temporal design:}
\begin{itemize}[leftmargin=1.5em]
    \item Student(student\_id, entity\_valid\_time)
    \item Instructor(instructor\_id, entity\_valid\_time)
    \item Advisor(student\_id, instructor\_id, valid\_time)
\end{itemize}

Temporal attributes may be represented separately, such as:
\begin{itemize}[leftmargin=1.5em]
    \item Student\_Name(student\_id, name, valid\_time)
    \item Student\_TotCred(student\_id, tot\_cred, valid\_time)
    \item Instructor\_Name(instructor\_id, name, valid\_time)
    \item Instructor\_Salary(instructor\_id, salary, valid\_time)
\end{itemize}

\textbf{Reduction to relations:}
\begin{itemize}[leftmargin=1.5em]
    \item Student(student\_id, entity\_valid\_from, entity\_valid\_to)
    \item Instructor(instructor\_id, entity\_valid\_from, entity\_valid\_to)
    \item Student\_Name(student\_id, valid\_from, name, valid\_to)
    \item Student\_TotCred(student\_id, valid\_from, tot\_cred, valid\_to)
    \item Instructor\_Name(instructor\_id, valid\_from, name, valid\_to)
    \item Instructor\_Salary(instructor\_id, valid\_from, salary, valid\_to)
    \item Advisor(student\_id, instructor\_id, valid\_from, valid\_to)
\end{itemize}

\textbf{Question 8}\\
Analyze, solve, and present the solution of Problem 6.14.

\textbf{Answer:}

\textbf{Superkey:} any set of attributes that uniquely identifies a tuple.

\textbf{Candidate key:} a minimal superkey.

\textbf{Primary key:} the candidate key selected as the main identifier.

\textbf{Example:}  
If Student(ID, email, name) has both ID and email unique, then:
\begin{itemize}[leftmargin=1.5em]
    \item ID is a candidate key
    \item email is also a candidate key
    \item (ID, email) is a superkey but not a candidate key
\end{itemize}

If ID is chosen as the main identifier, then it is the primary key.

\textbf{Question 9}\\
Analyze, solve, and present the solution of Problem 6.15.

\textbf{Answer:}

\insertdiagram{q6_15.png}{E-R diagram for Problem 6.15}

The hospital database may be modeled using:
\begin{itemize}[leftmargin=1.5em]
    \item Patient (\underline{patient\_id}, name, dob, sex, address, phone, blood\_group)
    \item Doctor (\underline{doctor\_id}, name, specialization, phone, department)
    \item Test\_Exam (\underline{log\_id}, test\_type, test\_date, result\_summary, notes)
\end{itemize}

Relationships:
\begin{itemize}[leftmargin=1.5em]
    \item For\_Patient(Test\_Exam, Patient)
    \item Conducted\_By(Test\_Exam, Doctor)
\end{itemize}

This model is clean because each test or examination is treated as an entity.

\textbf{Question 10}\\
Analyze, solve, and present the solution of Problem 6.19.

\textbf{Answer:}

Weak entities remain useful even when they can be transformed into strong entities because they naturally express:
\begin{itemize}[leftmargin=1.5em]
    \item existence dependence
    \item identification relative to an owner entity
    \item better conceptual clarity
    \item avoidance of artificial keys
\end{itemize}

Therefore, weak entities are important for semantic modeling.

\newpage
\textbf{Question 11}\\
Analyze, solve, and present the solution of Problem 6.21.

\textbf{Answer:}

\textbf{(a) Add Blu-ray discs and downloadable video}

\insertdiagram{q6_21a.png}{Extended E-R diagram for Problem 6.21(a)}

A suitable extension is:
\begin{itemize}[leftmargin=1.5em]
    \item Video\_Work(video\_id, title, year, category)
    \item BluRay\_Video(video\_id, price, region\_code)
    \item Download\_Video(video\_id, price, file\_format, drm\_type)
\end{itemize}

This allows the same content to exist in one or both formats with separate prices.

\textbf{(b) Allow shopping basket to contain books and video items}

\insertdiagram{q6_21b.png}{Extended E-R diagram for Problem 6.21(b)}

Introduce:
\begin{itemize}[leftmargin=1.5em]
    \item Store\_Item(item\_id, title)
\end{itemize}

Subclasses:
\begin{itemize}[leftmargin=1.5em]
    \item Book(item\_id, ISBN, year, price)
    \item BluRay(item\_id, price)
    \item DownloadVideo(item\_id, price)
\end{itemize}

Then replace:
contains(shopping\_basket, book)

with:
contains(shopping\_basket, store\_item, quantity)

\textbf{Question 12}\\
Analyze, solve, and present the solution of Problem 6.22 (ER diagram only).

\textbf{Answer:}

\insertdiagram{q6_22.png}{E-R diagram for Problem 6.22}

A suitable design includes:
\begin{itemize}[leftmargin=1.5em]
    \item Brand(brand\_id, brand\_name)
    \item Model(model\_id, brand\_id, model\_name, model\_year, body\_style, base\_price)
    \item Option(option\_id, option\_name, description, option\_price)
    \item Vehicle(VIN, model\_id, color, manufacture\_date, status)
    \item Dealer(dealer\_id, name, address, phone)
    \item Sales\_Staff(staff\_id, dealer\_id, name, phone)
    \item Customer(customer\_id, name, address, phone, email)
    \item Sales\_Order(order\_id, customer\_id, dealer\_id, staff\_id, VIN, order\_date, agreed\_price, order\_status)
\end{itemize}

Relationships:
\begin{itemize}[leftmargin=1.5em]
    \item Brand has many Models
    \item Model offers many Options
    \item Vehicle is an instance of one Model
    \item Vehicle may have many selected Options
    \item Dealer inventories Vehicles
    \item Dealer employs Sales\_Staff
    \item Customer places Sales\_Orders
    \item Sales\_Staff handles Sales\_Orders
\end{itemize}

\textbf{Question 13}\\
Analyze, solve, and present the solution of Problem 6.23 (ER diagram only).

\textbf{Answer:}

\insertdiagram{q6_23.png}{E-R diagram for Problem 6.23}

A suitable design includes:
\begin{itemize}[leftmargin=1.5em]
    \item Customer(customer\_id, name, address, phone, email)
    \item Shipment(shipment\_id, sender\_id, receiver\_id, ship\_date, service\_type)
    \item Package(tracking\_no, shipment\_id, weight, dimensions, status)
    \item Location(location\_id, location\_type, description)
    \item Truck(location\_id, registration\_no, route\_name)
    \item Plane(location\_id, tail\_no, airline\_code)
    \item Airport(location\_id, airport\_code, city, country)
    \item Warehouse(location\_id, warehouse\_code, address)
    \item Package\_Tracking(tracking\_event\_id, tracking\_no, location\_id, event\_time, event\_status)
\end{itemize}

Key idea: the current location of a package is the most recent tracking record.

\textbf{Question 14}\\
Analyze, solve, and present the solution of Problem 6.24 (ER diagram only).

\textbf{Answer:}

\insertdiagram{q6_24.png}{E-R diagram for Problem 6.24}

A suitable airline database design includes:
\begin{itemize}[leftmargin=1.5em]
    \item Airport(airport\_code, name, city, country)
    \item Aircraft(aircraft\_id, model, seat\_capacity)
    \item Flight\_Route(route\_id, flight\_no, origin\_code, destination\_code, scheduled\_duration)
    \item Flight\_Instance(instance\_id, route\_id, aircraft\_id, departure\_time, arrival\_time, status)
    \item Passenger(passenger\_id, name, phone, email, passport\_no)
    \item Booking(booking\_id, passenger\_id, instance\_id, seat\_no, booking\_class, booking\_status)
    \item Crew\_Member(crew\_id, name, role)
    \item Crew\_Assignment(crew\_id, instance\_id, duty)
\end{itemize}

Relationships:
\begin{itemize}[leftmargin=1.5em]
    \item A flight route has an origin airport and a destination airport
    \item A flight instance is a scheduled occurrence of a flight route
    \item An aircraft is assigned to a flight instance
    \item A passenger books a flight instance
    \item Crew members are assigned to flight instances
\end{itemize}

\newpage
\textbf{Question 15}\\
Analyze, solve, and present the solution of Problem 6.26.

\textbf{Answer:}

\insertdiagram{q6_26.png}{Specialization hierarchy for Problem 6.26}

A suitable generalization/specialization hierarchy is:

Motor\_Vehicle $\rightarrow$ Motorcycle, Passenger\_Car, Van, Bus

\textbf{Superclass attributes:}
\begin{itemize}[leftmargin=1.5em]
    \item VIN
    \item brand
    \item model
    \item model\_year
    \item color
    \item engine\_no
    \item fuel\_type
    \item price
\end{itemize}

\textbf{Subclass-specific attributes:}

\textbf{Motorcycle:}
\begin{itemize}[leftmargin=1.5em]
    \item engine\_capacity\_cc
    \item bike\_type
    \item has\_sidecar
\end{itemize}

\textbf{Passenger\_Car:}
\begin{itemize}[leftmargin=1.5em]
    \item number\_of\_doors
    \item trunk\_capacity
    \item body\_style
\end{itemize}

\textbf{Van:}
\begin{itemize}[leftmargin=1.5em]
    \item cargo\_volume
    \item payload\_capacity
    \item sliding\_door\_flag
\end{itemize}

\textbf{Bus:}
\begin{itemize}[leftmargin=1.5em]
    \item seating\_capacity
    \item standing\_capacity
    \item route\_service\_type
\end{itemize}

This placement is appropriate because common attributes stay in the superclass while specialized attributes remain in the subclasses.

\end{document}